\documentclass[../main.tex]{subfiles}
\graphicspath{{\subfix{../images/}}}
\begin{document}

\section{Nutrition and Health}

The institute of health promotion Switzerland\footnote{We are taking here Switzerland as an example for the health data.
  The original material has been compiled in Switzerland by Kilian Schmidt.
  Another point is in it's function as a representative western, first world country
  which is pretty healthy, and not so much affected by the first world problems as the USA.} leads on a regular base representative surveys
in the general population of Switzerland to the topic ``health and movement'',
in order to find out how sensibilized the general population is with respect to these topics.

The most important outcomes are:
\begin{itemize}
\item A big part of the swiss society seems to know a lot about the topics nutrition, movement and weight.
\item Two thirds of the people surveyed estimate that they can jusge for themselves, what is healthy and what is unhealthy nutrition.
\item Hald of the surveyed population says that they are informed about mevement for the sake of health benefits.
\item In average 80\% say that they can estimate themselves, if their weight is healthy.
\item Only one in three sees a need to act in their own case.
\end{itemize}

    \begin{figure}[htb!]
    \begin{center}
  \includegraphics[width=4cm]{MistakesNutrition}
  \caption{Our most important mistakes with food.}
  \end{center}
  \end{figure}


  \subsubsection{Nevertheless, the Swiss get more and more thick and sick.}

  \begin{itemize}
  \item The health survey of the Swiss Federal Institute for Statistik in the years 92, 97 and 2002
    showed a rapid increase in overweight (BMI 25--30) and obesity (adiposity) (BMI $>$ 30).
  \item In 2002 were 29.4\% of the Swiss adults overweight and 7.7\% obese.
  \item 2008, already every 5th child was overweight and every 20ieth child was obese.
    \item In the year 2000, 158,000 Swiss citiziens (2.5\% of the population) suffered from Diabetes type 2. This number doubles about every ten years.
  \end{itemize}

\end{document}